\section{Kaluza--Klein spectra, boundary conditions, and secular determinants}
\label{sec:kk-boundary}
% Status: cited/open. KK and interval background are cited; interval reductions remain theorem targets.

Kaluza--Klein theory supplies the second concrete language for the DeVries program.  The original idea unifies gauge fields with geometry by expanding a higher-dimensional theory around a compact ground state.  In Witten's account, symmetries of the compact space appear as gauge symmetries in the lower-dimensional theory, and the physical spectrum is obtained from small oscillations around the compactification background \cite{WittenKK1981}.  Salam and Strathdee give the complementary harmonic-expansion language on compact homogeneous spaces \(G/H\), where zero modes give gravitational plus Yang--Mills fields and higher modes give massive multiplets \cite{SalamStrathdee1982}.  The same account also records the hard model-building burden: obtaining the Standard Model gauge group, fermion quantum numbers, and scalar sector is a serious representation problem, even when the gauge group itself can be engineered.

The DeVries problem asks for a narrower and sharper object inside this tradition: a compact or interval derivation of the two-channel determinant
\begin{equation}
  \det
  \begin{pmatrix}
    \lambda & -\sqrt J\\
    -\sqrt J & \lambda+J
  \end{pmatrix}
  =
  \lambda^2+J\lambda-J.
  \label{eq:boundary-devries-target}
\end{equation}
This determinant is the entire target of the KK route.

\subsection{Circle compactification and two charge channels}
\label{sec:kk-circle-channels}

String compactification on a circle gives a clean source-backed example of a two-channel spectral structure.  Ordinary Kaluza--Klein momentum is quantized along the circle.  Closed strings add winding around the same circle.  Polchinski's circle example shows that the Kaluza--Klein vector couples to compact momentum, while the vector descending from the antisymmetric tensor couples to winding \cite{Polchinski1994}.  At the self-dual radius the spectrum has enhanced current algebra, and changing the radius is described by an operator built from the same currents.

This gives a source-backed analogy:
\begin{equation}
  \left(
  \begin{array}{c}
  \hbox{momentum-like channel}\\
  \hbox{winding-like or boundary channel}
  \end{array}
  \right)
  \quad\leadsto\quad
  \{x_+(J),x_-(J)\}.
  \label{eq:momentum-winding-channel}
\end{equation}
The analogy becomes a derivation after the two channels are reduced to Eq.~\eqref{eq:boundary-devries-target}.  The self-dual-radius example also explains why a scalar modulus belongs in the discussion.  The compact radius is a scalar field in the lower-dimensional theory, and gauge enhancement can occur at special values of that scalar.  This is the compactification version of the manuscript's positive/negative branch problem: a vector spectral statement and a scalar modulus can arise from the same compact sector.

A sharper test combines this channel picture with a boundary Schur complement.
Let the compact or interval sector contain a momentum/current-like channel and a
winding, brane-separation, or order-parameter channel.  After the tower is
integrated out, the reduced boundary kernel should have the form
\begin{equation}
  K^{\rm eff}_{j}(\lambda)
  =
  K^{\rm bdry}_{j}(\lambda)
  -
  V_j^\dagger(\lambda-L_j)^{-1}V_j
  \longrightarrow
  \begin{pmatrix}
  \lambda&-\sqrt{A_j}\\
  -\sqrt{A_j}&\lambda+A_j
  \end{pmatrix},
  \qquad A_j=j(j+1).
  \label{eq:kk-regge-schur-target}
\end{equation}
This target is the KK form of the Regge-intercept reading in
Sec.~\ref{sec:dual-model-lineage}: the compact or boundary calculation supplies
the two-channel intercept, and the oscillator or tower label supplies the
large-spin growth.

\subsection{Interval variational data}
\label{sec:interval-variational-data}

Extra-dimensional field theory on an interval gives a more direct determinant language.  Csaki--Hubisz--Meade derive boundary conditions from the variational principle for fields on an interval and show how boundary mass terms and boundary kinetic terms modify those conditions \cite{CsakiHubiszMeade2005}.  Dirichlet-Higgs and radion-stabilization models give adjacent source variables for boundary scalar profiles, radion data, and bulk Higgs stabilizers \cite{Bucci2004,HabaOda2011}.  A scalar example has the schematic action
\begin{equation}
  S_{\rm int}
  =
  \int d^4x\int_0^Ldy\,
  \frac12\Big[
  \partial_M\Phi\,\partial^M\Phi
  -V(\Phi)
  \Big]
  +S_0[\Phi]+S_L[\Phi],
  \label{eq:interval-action}
\end{equation}
with boundary functionals \(S_0\) and \(S_L\).  The boundary variation supplies equations at \(y=0\) and \(y=L\).  Boundary mass terms give Robin-type conditions.  Boundary kinetic terms can make the boundary equation depend on the four-dimensional eigenvalue.  This is exactly the structural feature needed by a secular determinant.

For a mode \(\Phi(x,y)=\phi_\lambda(x)f_\lambda(y)\), the interval problem has the abstract form
\begin{align}
  \mathcal L_J f_\lambda(y)&=\lambda f_\lambda(y),
  \qquad 0<y<L,
  \label{eq:interval-bulk-eigenproblem}\\
  \mathcal B_{0,J}(\lambda)f_\lambda(0)&=0,
  \qquad
  \mathcal B_{L,J}(\lambda)f_\lambda(L)=0.
  \label{eq:interval-boundary-eigenproblem}
\end{align}
The DeVries route asks for a two-dimensional light sector \(\mathcal H_J\) extracted from this problem.  A convenient basis is
\begin{equation}
  \mathcal H_J=\operatorname{span}\{h_J,a_J\},
  \label{eq:interval-light-sector}
\end{equation}
where \(h_J\) is a boundary order-parameter amplitude and \(a_J\) is an adjoint/current amplitude.  The interval derivation then reduces the bulk-plus-boundary problem to an effective kernel
\begin{equation}
  K_J^{\rm int}(\lambda)=
  \begin{pmatrix}
    \lambda+\Sigma_{hh,J}^{\rm int}(\lambda)
    &
    \Sigma_{ha,J}^{\rm int}(\lambda)
    \\
    \Sigma_{ah,J}^{\rm int}(\lambda)
    &
    \lambda+\Sigma_{aa,J}^{\rm int}(\lambda)
  \end{pmatrix}.
  \label{eq:interval-effective-kernel}
\end{equation}
The DeVries target is the boundary limit
\begin{equation}
  \Sigma_{hh,J}^{\rm int}=0,\qquad
  \Sigma_{aa,J}^{\rm int}=J,\qquad
  \Sigma_{ha,J}^{\rm int}\Sigma_{ah,J}^{\rm int}=J.
  \label{eq:interval-self-energy-target}
\end{equation}
This is the interval version of the endpoint target in Sec.~\ref{sec:endpoint-higgsing}.

The CHM source gives concrete boundary data that can feed this target.  In a
scalar toy problem, a boundary kinetic term produces an eigenvalue-dependent
condition and a modified product of the schematic form
\begin{equation}
  \partial_y f_\lambda\big|_{y_i}
  =
  -\frac{\lambda}{M_i}f_\lambda(y_i),
  \qquad
  (f,g)_{\rm CHM}
  =
  \int_0^L f(y)g(y)\,dy
  +
  \sum_i \frac{1}{M_i}f(y_i)g(y_i).
  \label{eq:chm-boundary-kinetic-data}
\end{equation}
For gauge fields, boundary scalar vevs give vector boundary data of the
schematic form
\begin{equation}
  \partial_y A_\mu\big|_{y_i}\mp v_i^2A_\mu(y_i)=0,
  \qquad
  \left(\frac{m^2}{\xi_i}-v_i^2\right)\pi_i+v_iA_5\big|_{y_i}=0,
  \label{eq:chm-boundary-vector-scalar-data}
\end{equation}
with endpoint-orientation signs fixed by the interval convention.  These
source facts define an allowable class of boundary kernels.  The scalar
equation also makes the gauge-fixing burden explicit: the projected DeVries
kernel must state the treatment of \(\xi_i\), separate physical scalar data
from eaten Goldstone data, preserve the photon zero mode, and map the interval
eigenvalue into the four-dimensional transverse pole condition.  The equality
to Eq.~\eqref{eq:interval-self-energy-target} is the open DeVries entry
theorem.

A Hodge factorization subroute gives a sharper test for the off-diagonal entry.
Witten's supersymmetric quantum mechanics on forms uses the exterior derivative
and its adjoint to build supercharges whose Hamiltonian is the Hodge Laplacian;
the Witten deformation adds Morse-function data to the same complex
\cite{WittenMorse1982}.  If an interval or compact sector supplies normalized
states \(e^0_J\in\Omega^0\) and \(e^1_J\in\Omega^1\) with
\begin{equation}
  D_Je^0_J=\sqrt J\,e^1_J,\qquad
  D_J^\dagger e^1_J=\sqrt J\,e^0_J,
  \label{eq:interval-hodge-factorization}
\end{equation}
then the de Rham part of the reduced operator supplies
\(\Sigma_{ha,J}^{\rm int}\Sigma_{ah,J}^{\rm int}=J\).  In Target X language,
Eq.~\eqref{eq:interval-hodge-factorization} supplies \(Q_{{\rm dR},J}\); the
boundary or breaking datum \(B_J\) must supply the diagonal entry \(-J\).  The
same interval model must derive the sign convention in the DeVries block and
the ordered electroweak map.  Appendix~\ref{app:target-susyqm-origin} records
this as a theorem target.

This gives the interval route an entry order.  The Hodge/SUSY-QM calculation
tests the square-root product
\(\Sigma_{ha,J}^{\rm int}\Sigma_{ah,J}^{\rm int}=J\).  The CHM boundary
calculation tests the diagonal current entry
\(\Sigma_{aa,J}^{\rm int}=J\).  A successful route must place both tests in
one source Hilbert space, with one inner product, one projection to
\(\mathcal H_J\), and one pole-matching map.  The negative branch is then read
from the eigenbasis of the same reduced kernel.

\subsection{Dimensional Schur-complement target}
\label{sec:kk-schur-target}

The dimensional interpolation of Appendix~\ref{app:dimensional-interpolation}
defines a candidate Schur-complement calculation.  Let the middle-dimensional sector be
split into a light boundary subspace \(\mathcal H_J=\operatorname{span}\{h_J,a_J\}\)
and a tower, bulk, or compact sector with operator \(L_J\).  With a boundary
kernel \(K^{\rm bdry}_J(\lambda)\), coupling map \(V_J\), and specified inner
product, integrating out the heavy sector gives
\begin{equation}
  K^{\rm eff}_J(\lambda)
  =
  K^{\rm bdry}_J(\lambda)
  -
  V_J^\dagger(\lambda-L_J)^{-1}V_J .
  \label{eq:interval-schur-complement}
\end{equation}
The required reduced kernel is
\begin{equation}
  K^{\rm eff}_J(\lambda)
  \longrightarrow
  \begin{pmatrix}
  \lambda&-\sqrt J\\
  -\sqrt J&\lambda+J
  \end{pmatrix}
  \label{eq:interval-schur-devries-target}
\end{equation}
in a limit where additional light channels decouple or acquire a controlled
block-diagonal form.  This formulation connects the \(D=10/6\) middle line to
the route kernel: \(K^{\rm bdry}_J\) contains boundary kinetic and mass data,
\(L_J\) contains the integrated tower, and \(V_J\) contains the
order-parameter/current mixing.  A successful interval derivation must specify
these three objects, the inner product, and the normalization yielding
Eq.~\eqref{eq:interval-self-energy-target}.

The same target can be stated in Dirichlet-to-Neumann form.  Solving the bulk
mode equation with fixed boundary values gives a map from boundary traces to
normal derivatives.  Boundary kinetic and scalar terms add local endpoint
pieces, so a source-level interval target is
\begin{equation}
  K^{\rm int}_J(\lambda)
  =
  P_J^\dagger
  \left(
  K^{\rm DtN}_J(\lambda)+K^{\rm brane}_J(\lambda)
  \right)P_J
  \longrightarrow
  \begin{pmatrix}
  \lambda&-\sqrt J\\
  -\sqrt J&\lambda+J
  \end{pmatrix},
  \label{eq:interval-dtn-target}
\end{equation}
where \(P_J\) is the projection to
\(\mathcal H_J=\operatorname{span}\{h_J,a_J\}\).  This form makes the
source-data demand explicit: the Dirichlet-to-Neumann kernel, brane terms,
projection, inner product, and decoupling rule must be supplied by one
interval model.

The current-channel entry can be isolated after the physical boundary kernel
has been converted to the dimensionless DeVries normalization.  Let \(P_{a,J}\)
select a transverse boundary gauge trace
\(a_J=P_{a,J}A_\mu^{T,\partial}\), normalized by the CHM product,
\(\langle a_J,a_J\rangle_{\rm CHM}=1\).  Let \(\Lambda_J\) denote the source
scale that makes the spectral variable dimensionless:
\begin{equation}
  \widehat\lambda=\frac{\lambda}{\Lambda_J^2}.
  \label{eq:chm-current-hat-lambda}
\end{equation}
With endpoint-normal signs and dimensionful brane terms fixed, define
\begin{equation}
  \widehat K^{\rm cur}_J(\widehat\lambda)
  =
  \frac{1}{\Lambda_J^2}
  \left\langle a_J,
  \left(
  K^{\rm DtN}_{T,J}(\lambda)
  +K^{\rm brane}_{T,J}(\lambda)
  -K^{\rm ref}_{T,J}(\lambda)
  \right)a_J
  \right\rangle_{\rm CHM},
  \qquad
  \lambda=\Lambda_J^2\widehat\lambda,
  \label{eq:chm-current-kernel}
\end{equation}
where \(K^{\rm ref}_{T,J}\) fixes the reference zero of the current kernel,
for example the unbroken photon channel or a stated boundary subtraction.
The source-supported CHM statement is the existence of this transverse
boundary kernel from the action, gauge fixing, boundary variation, and
modified product.  The DeVries-specific theorem target is
\begin{equation}
  \widehat K^{\rm cur}_J(\widehat\lambda)
  =
  \widehat\lambda+\widehat\Sigma_{aa,J}^{\rm CHM}
  +O(\widehat\lambda^2),
  \qquad
  \widehat\Sigma_{aa,J}^{\rm CHM}=J .
  \label{eq:chm-current-entry-normalized}
\end{equation}
The data still to be derived are \(P_{a,J}\), \(K^{\rm ref}_{T,J}\), the
source scale \(\Lambda_J\), endpoint-normal conventions, the photon projection,
the ordered W/Z boundary map, and the map
\(\widehat\lambda\mapsto m_n^2/\Lambda_J^2\mapsto\Pi_T^{(4)}(s)\).

\subsection{Dictionary for the interval entries}
\label{sec:interval-entry-dictionary}

The interval construction separates the entries of Eq.~\eqref{eq:boundary-devries-target} into physically addressable data:
\begin{center}
\begin{tabular}{p{0.24\linewidth}p{0.34\linewidth}p{0.32\linewidth}}
\toprule
Entry & Candidate source datum & DeVries-specific obligation \\
\midrule
\(\Sigma_{hh,J}^{\rm int}=0\) & boundary scalar reference or subtraction datum & fix the reference by a symmetry, a brane condition, or a normalization convention \\
\(\Sigma_{aa,J}^{\rm int}=J\) & boundary kinetic product, vector Robin term, or adjoint-current boundary datum & identify the quadratic invariant that appears in the boundary equation \\
\(\Sigma_{ha,J}^{\rm int}\Sigma_{ah,J}^{\rm int}=J\) & scalar-vector mixing, endpoint/bulk overlap, or brane-localized coupling & derive the invariant from a normalized overlap or generator product \\
\((J_H,J_{\rm adj})=(3/4,2)\) & Higgs/order-parameter channel and adjoint/current channel & derive the ordered electroweak sampling rule from the coupled interval problem \\
\bottomrule
\end{tabular}
\end{center}
The table is intentionally strict.  A boundary mass term by itself can produce many secular equations.  The DeVries claim requires the trace and determinant data in Eq.~\eqref{eq:interval-self-energy-target} and the electroweak ordering of \(J\)-samples.

The possible origins of \(J\) are correspondingly limited:
\begin{equation}
  J
  \in
  \left\{
  C_2(R_{\rm int}),
  C_2(R_{\rm bdry}),
  h_{\rm CFT},
  q_{\rm KK}^2,
  \langle f_{\rm bdry},f_{\rm bulk}\rangle^2
  \right\}.
  \label{eq:J-origin-list}
\end{equation}
Here \(C_2\) denotes a quadratic invariant, \(h_{\rm CFT}\) a conformal weight or current-algebra datum, \(q_{\rm KK}\) a compact charge, and the overlap term a normalized boundary/bulk integral.  A derivation names one of these entries and derives the equality to \(J\).

\subsection{Gauge fields, boundary breaking, and the electroweak ray}
\label{sec:kk-gauge-ray}

Gauge fields on an interval carry a stronger constraint than scalar toy models.  Boundary conditions that break gauge symmetry have to arise from the action and preserve the consistency of the lower-dimensional gauge theory.  Csaki--Hubisz--Meade use this language for electroweak symmetry breaking by boundary conditions and for Higgsless models \cite{CsakiHubiszMeade2005}.  Their framework is relevant here because it gives a controlled setting in which massive vector modes, a photon zero mode, custodial structure, and brane-localized terms can be discussed in one spectral problem.

The DeVries construction requires the boundary model to preserve the electroweak ray.  The low-energy gauge-Higgs mass matrix has the form
\begin{equation}
  \mathcal M_{\rm EW}^2(v)
  =
  v^2\,\mathcal P(g,g'),
  \label{eq:ew-ray-boundary}
\end{equation}
where \(v\) is radial and \(\mathcal P(g,g')\) is projective.  A completed
interval theorem would produce a boundary operator whose pole condition reduces
to this projective datum.  Thus the interval model is required to yield a
coupled charged/neutral construction:
\begin{equation}
  \left[
  \begin{array}{c}
  SU(2)_L\hbox{ charged sector}\\
  SU(2)_L\times U(1)_Y\hbox{ neutral sector}\\
  \hbox{photon zero mode}
  \end{array}
  \right]
  \quad\longrightarrow\quad
  \frac{M_W^2}{M_Z^2}
  =
  \frac{x_+(J_H)}{x_+(J_{\rm adj})}.
  \label{eq:interval-ew-target}
\end{equation}
This requirement places the W and Z readings inside one gauge-Higgs or boundary system.  The photon must remain the unbroken mode, and the massive neutral pole must carry the adjoint/current sample.

The corresponding proof obligation is
\begin{equation}
  S_{\rm 5D}+S_{\rm bdry}
  \longrightarrow
  \{
  \mathcal B_W(\lambda),
  \mathcal B_Z(\lambda),
  \mathcal B_\gamma(\lambda)
  \}
  \longrightarrow
  M_\gamma=0,\qquad
  \frac{M_W^2}{M_Z^2}
  =
  \frac{x_+(3/4)}{x_+(2)}.
  \label{eq:interval-ew-boundary-proof}
\end{equation}
The proof must include the charged and neutral boundary equations, the photon
zero mode, custodial conditions, the reduced inner product, and the pole
matching chain.

The interval version of the ordered assignment theorem is:
\begin{equation}
\begin{gathered}
  \text{bulk gauge fields}
  +\text{boundary terms}
  +\text{order-parameter channel}
  \\
  \Longrightarrow
  K_J^{\rm int}(\lambda)
  \Longrightarrow
  (J_W,J_Z)=(3/4,2).
\end{gathered}
\label{eq:interval-assignment-theorem}
\end{equation}
This theorem is parallel to the endpoint theorem in Sec.~\ref{sec:endpoint-higgsing}.  The two routes can eventually compete: endpoint labels emphasize matrix algebra, while interval boundary conditions emphasize variational origin and pole equations.

\subsection{Scalar partner from interval data}
\label{sec:kk-scalar-partner}

The negative branch also has a KK address.  Witten's review notes that compactification naturally produces scalar degrees of freedom, including radius moduli and possible charged scalar modes tied to symmetry breaking questions \cite{WittenKK1981}.  Polchinski's circle compactification likewise identifies the compact radius with a lower-dimensional scalar and shows that changing the radius is tied to current algebra at enhanced-symmetry points \cite{Polchinski1994}.  Hosotani's compact gauge-field mechanism adds a Wilson-line scalar route \cite{Hosotani1983}.  In interval language, scalar data can appear as boundary values, Wilson-line-like data, brane-localized order parameters, or moduli of the compact geometry \cite{Bucci2004,HabaOda2011}.

The scalar-branch target is therefore
\begin{equation}
  u_-(J)\in\mathcal H_J
  \quad\longmapsto\quad
  \mathcal F_{\rm KK}
  \in
  \left\{
  R(x),\,
  H^\dagger H,\,
  A_5,\,
  \hbox{brane scalar},\,
  \hbox{localized modulus}
  \right\}.
  \label{eq:kk-negative-branch-map}
\end{equation}
Here \(u_-(J)\) is the negative-branch eigenvector of the two-channel kernel.  A derivation identifies \(\mathcal F_{\rm KK}\), states its gauge transformation properties, and explains its scheme relation to the low-energy Higgs/order-parameter variables.

\subsection{KK completion criterion}
\label{sec:kk-completion-criterion}

The KK route is fixed by the following completion chain:
\begin{equation}
\begin{aligned}
  \text{compact or interval source}
  &\longrightarrow
  \text{boundary variational problem}
  \longrightarrow
  K_J^{\rm int}(\lambda),
  \\
  K_J^{\rm int}(\lambda)
  &\longrightarrow
  \lambda^2+J\lambda-J
  \longrightarrow
  \{\sPole,\mathcal F_{\rm KK}\}.
\end{aligned}
\label{eq:kk-completion-chain}
\end{equation}
The first arrow is source-backed by the KK and interval literature.  The second arrow is the DeVries determinant derivation.  The last arrow is the pole-placement and scalar-branch assignment.  This chain gives the manuscript a concrete way to compare compactification, endpoint, and \(G_2\) routes.
